\documentclass[UTF8, 12pt, fontset=fandol, oneside]{ctexart}
\usepackage{researchnotes}

\title{四种常见平均数及其幂平均统一形式}
\author{}
\date{}

\begin{document}

\maketitle

平均数是把一组数压缩成一个代表值的方法. 在不同问题中，``代表''的含义并不相同：有时我们关心总量的均分，有时关心比例的连续增长，有时关心速度、效率等倒数型量，有时又希望较大的观测值得到更高的惩罚. 因此，平均数不是唯一的概念，而是一族根据问题结构选择的数值概括方式.

设
\[
  x_1,x_2,\cdots,x_n
\]
是一组实数. 若需要使用几何平均、调和平均等涉及乘积或倒数的平均数，通常还要求 $x_i>0$. 若带有权重，则设
\[
  w_i\ge0,\qquad \sum_{i=1}^n w_i=1.
\]
权重 $w_i$ 表示第 $i$ 个数据在平均过程中的相对重要性.

\section*{一、算术平均}

\noindent\textbf{定义}\quad 对实数 $x_1,x_2,\cdots,x_n$，其算术平均数定义为
\[
  A=\frac{x_1+x_2+\cdots+x_n}{n}.
\]
若带权重 $w_1,w_2,\cdots,w_n$，则加权算术平均为
\[
  A_w=\sum_{i=1}^n w_i x_i.
\]

算术平均的核心思想是\textbf{总量均分}. 若 $x_i$ 表示每个人拥有的物品数量、每次测量的误差、若干门课程的成绩等，把总和平均分给 $n$ 个对象，所得数值就是算术平均.

\noindent\textbf{基本性质}\quad 算术平均具有线性性质：
\[
  \frac{(ax_1+b)+\cdots+(ax_n+b)}{n}
  =aA+b.
\]
这说明对数据整体进行线性变换后，平均数也作同样的线性变换. 算术平均还满足
\[
  \sum_{i=1}^n (x_i-A)=0,
\]
即各数据相对平均数的偏差和为零.

\noindent\textbf{例}\quad 某同学连续三天记录从家到学校的通勤时间，分别为 $28,32,30$ 分钟，则这三天的平均通勤时间为
\[
  A=\frac{28+32+30}{3}=30\text{ 分钟}.
\]
这表示若把三天总共 $90$ 分钟的通勤时间平均分到每天，每天相当于 $30$ 分钟.

若一门课程的总评由平时作业、期中考试、期末考试三部分组成，权重依次为 $20\%$、$30\%$、$50\%$. 若对应成绩分别为 $80,90,86$，则总评成绩为
\[
  0.2\cdot80+0.3\cdot90+0.5\cdot86=86.
\]
这里期末考试所占权重最大，因此总评成绩会更多地受到期末成绩的影响.

\section*{二、几何平均}

\noindent\textbf{定义}\quad 对正数 $x_1,x_2,\cdots,x_n$，其几何平均数定义为
\[
  G=\sqrt[n]{x_1x_2\cdots x_n}.
\]
带权几何平均为
\[
  G_w=\prod_{i=1}^n x_i^{w_i}.
\]

几何平均的核心思想是\textbf{乘法结构的均分}. 它适合处理倍数、比例、增长率等按乘法累积的量. 例如，一个数量先变为原来的 $x_1$ 倍，再在新的结果上变为原来的 $x_2$ 倍，依此继续到第 $n$ 步，那么最终会变为原来的
\[
  x_1x_2\cdots x_n
\]
倍.

几何平均回答的问题是：若把每一步都改成同一个倍率 $G$，并且要求 $n$ 步后的总效果不变，即
\[
  G^n=x_1x_2\cdots x_n,
\]
那么这个等效倍率 $G$ 应该是多少？因此，它描述的是倍率的平均，而不是百分数的简单平均.

\noindent\textbf{对数视角}\quad 因为
\[
  \ln G=\frac{1}{n}\sum_{i=1}^n \ln x_i,
\]
所以几何平均也可以理解为：先用对数把相乘的倍率转化成相加的量，再求算术平均，最后指数还原. 这个视角不是改变定义，而是说明几何平均为什么适合描述连续按比例变化的过程.

\noindent\textbf{例}\quad 某投资第一年增长 $20\%$，第二年下降 $10\%$，第三年增长 $50\%$. 三年的增长倍率分别为
\[
  1.2,\quad 0.9,\quad 1.5.
\]
等效的年均增长倍率为
\[
  G=\sqrt[3]{1.2\cdot0.9\cdot1.5}=\sqrt[3]{1.62}.
\]
因此等效年均增长率为
\[
  \sqrt[3]{1.62}-1.
\]
这里不能直接取 $20\%,-10\%,50\%$ 的算术平均，因为本金在每一年都会被前一年的变化重新缩放.

\section*{三、调和平均}

\noindent\textbf{定义}\quad 对正数 $x_1,x_2,\cdots,x_n$，其调和平均数定义为
\[
  H=\frac{n}{\frac{1}{x_1}+\frac{1}{x_2}+\cdots+\frac{1}{x_n}}.
\]
带权调和平均为
\[
  H_w=\frac{1}{\sum_{i=1}^n \frac{w_i}{x_i}}.
\]

调和平均的核心思想是\textbf{倒数结构的均分}. 当数据本身是``单位资源产生多少结果''或``单位结果消耗多少资源''这类比率时，若固定的是分子或分母中的某一总量，调和平均往往是合适的平均方式.

\noindent\textbf{例}\quad 一辆车去程速度为 $60$ 千米/小时，返程速度为 $40$ 千米/小时，且去程和返程路程相同. 平均速度不是
\[
  \frac{60+40}{2}=50,
\]
而是
\[
  H=\frac{2}{\frac{1}{60}+\frac{1}{40}}=48.
\]
这是因为平均速度等于总路程除以总时间，而时间与速度的倒数成正比.

\section*{四、平方平均}

\noindent\textbf{定义}\quad 对实数 $x_1,x_2,\cdots,x_n$，其平方平均数，也称均方根，定义为
\[
  Q=\sqrt{\frac{x_1^2+x_2^2+\cdots+x_n^2}{n}}.
\]
带权平方平均为
\[
  Q_w=\sqrt{\sum_{i=1}^n w_i x_i^2}.
\]

平方平均先平方、再求算术平均、最后开平方：平方使正负号不再互相抵消，开方则把结果还原到原来的量纲. 因此，当处理偏差、误差、周期性震荡等含有正负号的数据时，平方平均比普通算术平均更能反映数据的典型大小、波动强度或整体偏离程度. 由于平方会放大较大数的影响，它对大值更敏感，常用于交流电有效值、误差大小、向量长度以及统计中的标准差相关计算.

\noindent\textbf{例}\quad 某工厂加工一批零件，设计长度为 $100$ 毫米. 抽查三个零件后发现，它们相对设计长度的误差分别为 $-2,0,2$ 毫米. 若只取误差的算术平均，则
\[
  \frac{-2+0+2}{3}=0,
\]
这只能说明这三个零件整体上没有明显偏长或偏短，并不表示加工没有误差. 若关心零件偏离设计长度的典型大小，则应避免正负误差互相抵消，可用平方平均：
\[
  Q=\sqrt{\frac{(-2)^2+0^2+2^2}{3}}=\sqrt{\frac{8}{3}}\approx1.63.
\]
这表示这组三个零件的典型长度误差约为 $1.63$ 毫米.

\section*{五、幂平均}

\noindent\textbf{定义}\quad 对正数 $x_1,x_2,\cdots,x_n$ 和实数 $p$，$p$ 阶幂平均定义为
\[
  M_p=
  \begin{cases}
    \left(\dfrac{x_1^p+x_2^p+\cdots+x_n^p}{n}\right)^{1/p}, & p\ne0,\\[1.2em]
    \sqrt[n]{x_1x_2\cdots x_n}, & p=0.
  \end{cases}
\]
带权形式为
\[
  M_p(w)=
  \begin{cases}
    \left(\displaystyle\sum_{i=1}^n w_i x_i^p\right)^{1/p}, & p\ne0,\\[1.2em]
    \displaystyle\prod_{i=1}^n x_i^{w_i}, & p=0.
  \end{cases}
\]

幂平均把许多常见平均数统一到同一个框架中：
\[
  M_1=A,\qquad M_0=G,\qquad M_{-1}=H,\qquad M_2=Q.
\]
当 $p$ 越大时，大的观测值影响越强；当 $p$ 越小时，小的观测值影响越强.

\noindent\textbf{幂平均不等式}\quad 若 $p<q$，则
\[
  M_p\le M_q.
\]
特别地，对正数 $x_1,\cdots,x_n$ 有
\[
  H\le G\le A\le Q.
\]
等号成立当且仅当
\[
  x_1=x_2=\cdots=x_n.
\]
这条链式不等式是比较各种平均数时最重要的基本事实之一.

\subsection*{幂平均不等式的证明}

\begingroup
\color{rnorange}

\noindent\textbf{证明}\quad 令 $y_i=\ln x_i$，并定义
\[
  \varphi(t)=\ln\left(\frac{e^{t y_1}+e^{t y_2}+\cdots+e^{t y_n}}{n}\right).
\]
若记
\[
  \alpha_i(t)=\frac{e^{t y_i}}{e^{t y_1}+e^{t y_2}+\cdots+e^{t y_n}},
\]
则
\[
  \varphi''(t)
  =\sum_{i=1}^n\alpha_i(t)y_i^2
  -\left(\sum_{i=1}^n\alpha_i(t)y_i\right)^2\ge0.
\]
因此 $\varphi$ 是凸函数. 又 $\varphi(0)=0$，当 $t\ne0$ 时，
\[
  \ln M_t
  =\frac{1}{t}\ln\left(\frac{x_1^t+x_2^t+\cdots+x_n^t}{n}\right)
  =\frac{\varphi(t)}{t}.
\]
当 $t=0$ 时，$\ln M_0=(\ln x_1+\cdots+\ln x_n)/n=\varphi'(0)$. 凸函数的割线斜率随端点向右移动而不减，所以 $t\mapsto \ln M_t$ 是不减函数. 若 $p<q$，则
\[
  \ln M_p\le \ln M_q.
\]
由于对数函数单调递增，故
\[
  M_p\le M_q.
\]
取 $p=-1,0,1,2$，便得到
\[
  H=M_{-1}\le M_0=G\le M_1=A\le M_2=Q.
\]
若 $x_1,\cdots,x_n$ 不全相等，则 $\varphi$ 严格凸，上述不等式为严格不等式；因此等号成立当且仅当 $x_1=x_2=\cdots=x_n$.

\endgroup

\end{document}
